Este capítulo revisa críticamente la literatura correspondiente a las seis líneas identificadas en la Fase 1 del proyecto de tesis: (1) el fundamento técnico de RAG, (2) los sistemas de recomendación y emparejamiento de empleo basados en NLP y LLMs, (3) los límites del emparejamiento por palabra clave y de los ATS en reclutamiento, (4) Design Science Research como paradigma metodológico, (5) las bases de datos vectoriales, y (6) la clasificación híbrida de consultas y el text-to-SQL. Cada sección cierra con un párrafo que vincula explícitamente los hallazgos revisados con una decisión de diseño concreta de JobAI.

\section{Fundamento técnico de RAG}

El paradigma de \textit{Retrieval-Augmented Generation} fue introducido por \citet{Lewis2020} como una arquitectura que combina un recuperador denso (\textit{Dense Passage Retriever}) con un modelo generativo tipo secuencia-a-secuencia (BART), entrenados de forma conjunta para producir respuestas fundamentadas en pasajes recuperados de una base de conocimiento no paramétrica —en su caso, Wikipedia. Los autores demuestran que este diseño supera a los modelos puramente paramétricos en tareas intensivas en conocimiento (\textit{knowledge-intensive NLP tasks}), y que la separación entre memoria paramétrica (el modelo) y memoria no paramétrica (el índice de recuperación) permite actualizar el conocimiento del sistema sin reentrenar el modelo generativo, simplemente actualizando el índice.

La encuesta de \citet{Gao2023} sistematiza la evolución posterior de este paradigma en tres categorías: RAG \textit{ingenuo}, que recupera y genera en un único paso sin refinamiento adicional; RAG \textit{avanzado}, que incorpora preprocesamiento de la consulta (reformulación, expansión), reordenamiento de los resultados recuperados y compresión o filtrado del contexto antes de pasarlo al modelo generativo; y RAG \textit{modular}, que descompone el pipeline en componentes intercambiables y permite flujos condicionales, como enrutar la consulta hacia distintas estrategias de recuperación según su naturaleza. Los autores identifican como retos abiertos del paradigma la calidad de la recuperación (fragmentos irrelevantes degradan la respuesta final incluso con un modelo generativo capaz), la gestión de la longitud de contexto, y la persistencia de alucinaciones cuando el contexto recuperado es insuficiente o contradictorio.

Evaluar un sistema RAG es, en sí mismo, un problema no trivial: a diferencia de tareas con una respuesta de referencia única, la calidad de una respuesta generada depende de si está \emph{fundamentada} en el contexto recuperado y de si el contexto recuperado era, a su vez, relevante para la consulta. \citet{Es2024} proponen RAGAS, un framework de evaluación sin referencia humana que opera a partir de tres dimensiones: fidelidad (\textit{faithfulness}, el grado en que cada afirmación de la respuesta puede inferirse del contexto recuperado), relevancia de la respuesta (\textit{answer relevancy}) y precisión/exhaustividad del contexto recuperado (\textit{context precision/recall}). Estas métricas son especialmente relevantes para este trabajo porque no requieren un conjunto de referencia anotado manualmente para cada consulta posible sobre un catálogo de ofertas laborales que cambia cada seis horas.

\textbf{Relación con JobAI.} El pipeline implementado en JobAI —descrito en detalle en el Capítulo 4— corresponde, en la taxonomía de \citet{Gao2023}, a un diseño de RAG modular: la consulta del usuario se aumenta con información del perfil (títulos y habilidades extraídos del currículo) antes de embeberse, el contexto recuperado se somete a una etapa de detección de campos que complementa fragmentos incompletos con datos estructurados (salario, beneficios, requisitos), y la consulta completa se enruta condicionalmente entre una ruta de recuperación semántica y una ruta de agregación SQL (véase la Sección~\ref{sec:estado-arte-hibrido}). El diseño de las métricas de evaluación del Capítulo 3 —precisión de recuperación y fidelidad de las respuestas— se apoya directamente en las dimensiones de fidelidad y precisión de contexto propuestas por RAGAS \citep{Es2024}, adaptadas al dominio de ofertas laborales.

\section{Sistemas de recomendación y emparejamiento de empleo con NLP y LLMs}

La industria del reclutamiento en línea ha sido, durante la última década, uno de los espacios donde la búsqueda semántica basada en representaciones densas se ha adoptado con mayor decisión, precisamente por la limitación estructural del emparejamiento léxico descrita en la Sección~\ref{sec:estado-arte-keyword}. \citet{HaThuc2017} documentan la transición de LinkedIn Talent Search desde un modelo de \textit{Query-By-Keyword} —donde el reclutador construye manualmente una consulta booleana— hacia uno de \textit{Query-By-Example}, donde el reclutador simplemente indica uno o varios candidatos ideales y el sistema aprende a generar la consulta y el modelo de \textit{ranking} a partir de esos ejemplos. Los autores destacan como reto principal la ausencia de datos de uso históricos para un paradigma completamente nuevo, resuelto mediante una estrategia de aprendizaje que reutiliza registros del sistema anterior.

En una línea similar mas centrada en la representación semántica de vacantes, \citet{Bocharova2023} proponen VacancySBERT, una red siamesa basada en \textit{Sentence-BERT} entrenada para representar títulos de vacantes y habilidades en un espacio vectorial donde la similitud coseno refleja equivalencia semántica, incluso cuando el título exacto difiere ("Software Engineer" y "Backend Developer", por ejemplo). Los autores reportan mejoras de hasta 21,5\,\% sobre \textit{baselines} como FastText, BERT y SentenceBERT genéricos al incorporar las habilidades mencionadas en la oferta como característica adicional durante la inferencia, evidencia directa de que una representación especializada al dominio de reclutamiento supera a una representación general de propósito amplio.

A una escala mayor, \citet{Shen2024} describen cómo LinkedIn incorporó técnicas de recuperación aprendida (\textit{learning to retrieve}) —incluyendo recuperación densa y recuperación generativa— a su sistema de búsqueda y recomendación de empleo, en sustitución parcial del paradigma clásico de dos etapas (indexación invertida por entidades extraídas más re-ranking), que los autores señalan como costoso de mantener manualmente y limitado frente a consultas que no calzan exactamente con los términos indexados.

Finalmente, la revisión bibliométrica de \citet{RojasGaleano2022} sitúa estos desarrollos dentro de una tendencia más amplia: analizando la producción científica sobre inteligencia artificial aplicada al emparejamiento empleo-currículo, los autores documentan un crecimiento sostenido del área en la última década, concentrado mayoritariamente en tareas de análisis de currículos, extracción de habilidades y \textit{ranking}/recomendación mediante representaciones semánticas, con una presencia todavía incipiente de arquitecturas conversacionales o generativas completas aplicadas específicamente a este dominio.

\textbf{Relación con JobAI.} La estrategia de aumentar la consulta del usuario con los títulos y habilidades extraídos de su currículo antes de embeberla —descrita en el Capítulo 4— sigue el mismo principio que motiva a \citet{HaThuc2017} y \citet{Bocharova2023}: representar semánticamente tanto la necesidad del usuario como la oferta, en lugar de exigir una coincidencia textual entre ambas. La brecha señalada por \citet{RojasGaleano2022} respecto a la escasez de arquitecturas conversacionales completas en este dominio es, precisamente, el espacio que este trabajo busca documentar y evaluar empíricamente.

\section{Límites del emparejamiento por palabra clave y de los ATS en reclutamiento}
\label{sec:estado-arte-keyword}

Más allá de que la literatura de matching semántico (Sección anterior) motive sus propuestas contrastándolas con el emparejamiento por palabra clave, vale la pena precisar en qué consiste exactamente esa limitación. El problema, conocido en recuperación de información como \textit{desajuste de vocabulario} (\textit{vocabulary mismatch}), ocurre porque un sistema de coincidencia léxica exacta —o de expansión mediante listas fijas de sinónimos— no puede anticipar todas las formas en que dos personas describirán el mismo concepto. \citet{HaThuc2017} lo ilustran desde la perspectiva del reclutador: traducir un conjunto complejo de criterios de contratación a una consulta booleana efectiva es una tarea que requiere pericia y que, aun así, omite candidatos cuyo perfil describe la misma experiencia con palabras distintas a las de la consulta. \citet{Bocharova2023} lo cuantifican desde la perspectiva del catálogo de vacantes: dos títulos de cargo pueden ser funcionalmente equivalentes y, sin embargo, no compartir ningún término, lo que un ATS basado en coincidencia exacta interpreta simplemente como "no encontrado".

Esta limitación se traduce en dos tipos de error complementarios: falsos negativos (una oferta relevante no aparece porque no comparte el vocabulario exacto de la consulta) y falsos positivos de baja calidad (una oferta aparece por compartir una palabra clave aislada, sin que el resto de su contenido sea relevante). Ninguno de los dos es evidente para quien realiza la búsqueda, que simplemente percibe un sistema con baja precisión, sin distinguir si el problema es de recuperación insuficiente o de ausencia de contexto para razonar sobre la relevancia.

\textbf{Relación con JobAI.} Esta limitación es, textualmente, el punto de partida del proyecto (véase la Sección~\ref{sec:problema}): la premisa de que un panel de filtros y una caja de búsqueda por palabra clave no pueden expresar consultas del tipo \textit{"qué roles de backend en mi stack pagan más de \$120k"} o \textit{"cuáles de mis ofertas guardadas calzan con mi experiencia en Python"}, que requieren razonar sobre significado y no solo sobre coincidencia textual.

\section{Bases de datos vectoriales: de una extensión relacional a un motor dedicado}
\label{sec:estado-arte-vectordb}

Un sistema RAG requiere, en su componente de recuperación, resolver eficientemente el problema de búsqueda del vecino más cercano sobre potencialmente millones de vectores de \textit{embedding}. \citet{Malkov2020} proponen HNSW (\textit{Hierarchical Navigable Small World graphs}), un algoritmo de búsqueda aproximada del vecino más cercano basado en un grafo de proximidad organizado en capas jerárquicas, que permite una complejidad de búsqueda logarítmica y que se ha convertido en el índice de facto de la mayoría de motores vectoriales modernos, incluido Qdrant. \citet{Johnson2019}, con FAISS, abordan el mismo problema desde una perspectiva distinta —el uso de GPUs para paralelizar tanto la búsqueda exhaustiva como la búsqueda aproximada mediante cuantización de producto—, demostrando que es posible construir un grafo de vecinos más cercanos sobre mil millones de vectores en menos de doce horas usando hardware especializado. \citet{Aumuller2020}, con ANN-Benchmarks, aportan un marco de evaluación empírica estandarizado que permite comparar de forma reproducible el desempeño de estos algoritmos (y de las bibliotecas que los implementan) bajo distintos conjuntos de datos y regímenes de precisión-velocidad, una metodología que este trabajo adapta en el Capítulo 3 para la comparación cuantitativa entre pgvector y Qdrant.

Estos algoritmos pueden implementarse embebidos en un motor de base de datos relacional de propósito general —extendido con un tipo de dato vectorial, como hace pgvector sobre PostgreSQL— o en un motor construido específicamente para datos vectoriales. \citet{Wang2021} describen con Milvus el diseño de esta segunda familia: una arquitectura que separa el almacenamiento del cómputo, soporta múltiples tipos de índice ANN de forma nativa y está optimizada para cargas de trabajo heterogéneas (CPU y GPU) a una escala que, según los autores, resulta difícil de sostener con extensiones vectoriales sobre motores relacionales no diseñados originalmente para ese patrón de acceso. La encuesta de \citet{Pan2024} sistematiza más de veinte sistemas de este tipo (incluyendo Milvus y Qdrant), proponiendo una taxonomía según su modelo de indexación, su soporte de búsqueda híbrida (vectorial más filtros estructurados) y su estrategia de escalamiento horizontal; mientras que la revisión de \citet{Taipalus2024} ofrece una introducción conceptual a los fundamentos, casos de uso y retos actuales de estos sistemas, señalando la consistencia transaccional y la actualización incremental del índice como puntos todavía activos de investigación.

El trabajo más directamente relevante para justificar una decisión de migración es el de \citet{Sanca2024}, quienes estudian empíricamente las rutas de acceso a datos en escenarios de búsqueda mixta vector-relacional —es decir, consultas que combinan similitud vectorial con filtros sobre atributos estructurados, exactamente el patrón de acceso de una plataforma de empleo que necesita combinar "semánticamente similar a la consulta" con "salario mayor a X" o "modalidad remota"—. Los autores muestran que, a medida que crece el volumen de datos y la selectividad de los filtros estructurados varía, las rutas de acceso disponibles en motores relacionales de propósito general presentan compromisos de rendimiento que un motor construido específicamente en torno a índices ANN puede evitar por diseño, al tratar la búsqueda vectorial como ciudadano de primera clase en lugar de como una extensión añadida sobre un modelo de datos relacional preexistente.

\textbf{Relación con JobAI.} JobAI almacenó inicialmente sus \textit{embeddings} —de 768 dimensiones, generados para cada fragmento de oferta laboral y cada currículo— directamente en PostgreSQL mediante pgvector, sin un índice HNSW o IVFFlat configurado explícitamente, lo que en la práctica implica una búsqueda por barrido secuencial: viable a decenas de miles de fragmentos, pero un compromiso real de escalabilidad ante un catálogo que crece cada seis horas con cada ciclo de ingesta. La evidencia revisada en esta sección —particularmente el hallazgo de \citet{Sanca2024} sobre búsqueda mixta vector-relacional, y la arquitectura de referencia de \citet{Wang2021}— fundamenta la decisión, documentada en el Capítulo 3, de migrar el almacenamiento vectorial a Qdrant como iteración de diseño dentro del ciclo de Design Science Research, y de conservar temporalmente el backend de pgvector para ejecutar la comparación cuantitativa de latencia y \textit{recall} entre ambos que se reporta en el Capítulo 5, empleando la metodología de benchmarking de \citet{Aumuller2020} como referencia.

\section{Design Science Research como paradigma metodológico}

Design Science Research (DSR) se distingue de la investigación puramente conductual en sistemas de información en que su objetivo no es explicar o predecir un fenómeno organizacional, sino diseñar y evaluar rigurosamente un artefacto que resuelve un problema relevante. \citet{Hevner2004} formalizan este paradigma mediante siete lineamientos —entre ellos, que el artefacto de diseño debe ser una contribución identificable (un constructo, modelo, método o instanciación), que su utilidad y calidad deben demostrarse mediante métodos de evaluación rigurosos, y que la investigación debe comunicarse tanto a audiencias orientadas a la tecnología como a la gestión—, y proponen un marco conceptual de dos ciclos complementarios: el ciclo de \textit{relevancia}, que conecta el artefacto con las necesidades del entorno de aplicación, y el ciclo de \textit{rigor}, que lo conecta con la base de conocimiento científico existente.

\citet{Peffers2007} operacionalizan estos principios en una metodología de seis actividades secuenciales pero no estrictamente lineales: identificación del problema y motivación, definición de los objetivos de una solución, diseño y desarrollo del artefacto, demostración de su uso en un contexto adecuado, evaluación de qué tan bien el artefacto resuelve el problema, y comunicación de los resultados. Un aporte relevante de esta metodología para el presente trabajo es que reconoce explícitamente distintos puntos de entrada al ciclo —incluyendo un punto de entrada centrado en el artefacto (\textit{design and development centered}), apropiado cuando la investigación parte de un artefacto ya existente y busca evaluarlo o mejorarlo, en lugar de partir de un problema completamente nuevo—.

El volumen editado por \citet{HevnerChatterjee2010} consolida, cinco años después del artículo seminal de 2004, tanto la teoría como casos de aplicación práctica de DSR, reafirmando el marco de los tres ciclos (relevancia, diseño y rigor) como la forma estándar de comunicar cómo un proyecto de DSR se nutre de la literatura existente (rigor) y retorna a ella una contribución evaluada (relevancia), a la vez que documenta el ciclo interno de construcción y evaluación iterativa del artefacto mismo.

\textbf{Relación con JobAI.} Este trabajo adopta el punto de entrada centrado en el artefacto de \citet{Peffers2007}: JobAI ya existe como instanciación funcional (el artefacto, en el sentido de \citealp{Hevner2004}), construida originalmente sin un proceso de evaluación académica formal. Las actividades de diseño y desarrollo de este ciclo de DSR corresponden a la migración del almacenamiento vectorial a Qdrant (Sección~\ref{sec:estado-arte-vectordb}); las actividades de demostración y evaluación corresponden a la ejecución del experimento comparativo contra un \textit{baseline} de palabra clave y a la medición de las métricas de fidelidad y precisión de recuperación (Capítulo 3); y la comunicación de resultados corresponde al presente documento. El Capítulo 3 desarrolla con mayor detalle esta correspondencia entre las actividades del modelo de \citet{Peffers2007} y las fases del cronograma de este proyecto.

\section{Clasificación híbrida de consultas y text-to-SQL}
\label{sec:estado-arte-hibrido}

No todas las preguntas que un usuario formula sobre un catálogo de ofertas laborales son preguntas de recuperación semántica. Una pregunta como \textit{"¿cuántas ofertas remotas de Python hay en la base de datos?"} tiene una respuesta exacta y verificable mediante una agregación SQL, y forzarla por una ruta de recuperación vectorial —que por diseño devuelve los fragmentos más similares, no un conteo exacto— produciría una respuesta aproximada o inconsistente. \citet{Biswal2024} formalizan esta observación con TAG (\textit{Table-Augmented Generation}), un paradigma que argumenta que tanto el text-to-SQL puro como el RAG puro cubren, cada uno por separado, solo un subconjunto estrecho de las preguntas que un usuario real querría formular sobre una base de datos: el primero se limita a lo expresable en álgebra relacional exacta, y el segundo se limita a preguntas resolubles mediante la recuperación de uno o pocos registros puntuales. Los autores proponen unificar ambos caminos —consulta exacta sobre datos estructurados y síntesis semántica sobre datos no estructurados— bajo un mismo marco de razonamiento.

La dificultad práctica de resolver la mitad "exacta" de ese problema mediante un LLM sin restricciones queda documentada por \citet{Li2023} con BIRD, un banco de pruebas de text-to-SQL a gran escala (12.751 pares de pregunta-consulta sobre 95 bases de datos reales de 37 dominios profesionales). Los autores muestran que incluso modelos de lenguaje muy capaces obtienen un desempeño sustancialmente inferior al humano al tener que generar SQL ejecutable y correcto sobre esquemas de bases de datos reales —con sus inconsistencias, ambigüedades y volúmenes de datos genuinos—, en contraste con bancos de pruebas anteriores que operaban sobre esquemas simplificados con pocos registros. Esta brecha entre el text-to-SQL de laboratorio y el text-to-SQL sobre datos reales motiva preferir, en un sistema de producción, una clasificación de la intención de la consulta antes de decidir si se genera SQL, en lugar de delegar esa decisión íntegramente a un LLM de propósito general.

\textbf{Relación con JobAI.} El clasificador de consultas de JobAI —descrito en el Capítulo 4— resuelve este mismo problema mediante un diseño pragmático de dos etapas: un conjunto de patrones \textit{regex} captura las señales léxicas inequívocas de intención de agregación ("cuántos", "promedio", "más común") y enruta esas consultas directamente a una plantilla SQL, evitando por completo el riesgo de generación de SQL abierta que documenta \citet{Li2023}; solo cuando ningún patrón coincide, un clasificador basado en LLM decide entre recuperación, agregación o una ruta híbrida que combina ambas en un mismo prompt. Este diseño puede leerse como una instanciación deliberadamente conservadora del espíritu de TAG \citep{Biswal2024} —reconocer que recuperación semántica y consulta exacta son necesarias para cubrir el espacio completo de preguntas de un usuario—, que privilegia la determinación y la seguridad de un conjunto acotado de plantillas SQL sobre la generalidad, más frágil, de un LLM generando SQL libremente sobre el esquema completo de la base de datos.

\section{Síntesis del marco de referencia}

Las seis líneas revisadas en este capítulo convergen en la premisa que motiva este trabajo: la limitación del emparejamiento léxico en la búsqueda de empleo (Secciones~\ref{sec:estado-arte-keyword} y siguientes) tiene, en la literatura de recuperación de información y de RAG, una solución arquitectónica madura y ampliamente adoptada por la industria del reclutamiento a gran escala; esa solución requiere, a su vez, decisiones de infraestructura (almacenamiento vectorial) y de enrutamiento de consultas (clasificación híbrida) que la literatura reciente permite fundamentar y no solo intuir. Lo que la literatura revisada no ofrece todavía —y que constituye el aporte específico de este trabajo— es una evaluación cuantitativa integrada de estas piezas sobre un artefacto real y en producción, siguiendo un proceso metodológico explícito de Design Science Research. El Capítulo 3 traduce esta síntesis en un diseño metodológico concreto.
